\section{Ethics and Broader Impact}

\paragraph{Source content and consent.}
The fresh portion of the \method{} corpus was scraped from public
Facebook, TikTok, and YouTube comments on public pages. We removed
user handles and profile-photo URLs, replaced any quoted personal
name with a single token \texttt{[PER]}, and dropped any comment that
included a phone number, email, or address. The dataset is released
under CC-BY-NC-4.0 with a take-down clause; any user who recognises
their comment can request removal, and the released version is
regenerated monthly.

\paragraph{Annotator welfare.}
Three annotators (two reviewers and one adjudicator), all L1
Vietnamese, were paid $\$15$/hour and warned that the corpus contains
hate-speech and mocking content. After pre-filtering with ViHSD
\citep{vihsd2021} and ViHOS \citep{vihos2023} lists, the actual rate
of overtly toxic content is $3.4\%$; the skip rate is $1.1\%$.

\paragraph{Dual-use risk.}
A Vietnamese sarcasm and implicit-sentiment classifier could moderate
online discourse but also flag dissent under regimes that treat
irony as a workaround for direct criticism. The released model card
\citep{mitchell2019modelcards} recommends that \method{} not be used
as the sole input to enforcement actions on user accounts.

\paragraph{Bias and provenance.}
The corpus over-represents younger urban speakers and under-represents
older speakers and rural dialects \citep{vimd2024}. All datasets carry
an explicit licence in \texttt{data/manifest/datasets.json}; rationales
carry generator model ID, decoding parameters, and prompt hash. We
follow the data-statement and datasheet practice of
\citet{bender2018datastatements}, \citet{gebru2021datasheets}, and
\citet{bender2021stochastic}.
